\section{Related Work}
\label{sec:related}

\subsection{Wildlife Survey from Imagery}
\label{sec:rw:survey}

Every field method for estimating abundance carries assumptions whose violation biases the
result. Distance sampling fits a detection function to distances recorded along transects
and extrapolates to density. The \emph{Distance} software implements conventional,
multiple-covariate and mark--recapture engines, each relaxing more assumptions, but even the
most flexible still requires near-certain detection on the line, no responsive movement and
error-free distances~\cite{thomas2010distance}. Buckland et al.~\cite{buckland2023wildlife}
argue that camera, acoustic and spatial capture--recapture surveys are displacing classical
distance sampling faster than statistical ecology is adapting. For unmarked populations,
N-mixture models treat site-level population size as a latent Poisson variable estimated
from replicated counts~\cite{royle2004nmixture}. Spatial capture--recapture extends this to
camera-trap encounters but requires individually identifiable
animals~\cite{royle2009bayesian}, which thermal imagery at our resolution cannot supply.
Iijima~\cite{iijima2020review} reviews the estimators available for unmarked populations
and concludes that none is universally applicable. Two recent corrections apply directly
to road transects. Movement compensation corrects the detection inflation caused by animals
encountering a moving observer~\cite{morgan2024wildlife}. Animals may also be attracted to
or repelled by the road itself, which violates the uniform-distribution
assumption~\cite{diefenbach2025accounting}.

Thermal cameras detect animals by emitted heat and so allow nocturnal survey where visual
methods fail. Burke et al.~\cite{burke2019optimizing} identify the binding constraints:
contrast loss as ground temperature approaches body temperature, atmospheric absorption,
vegetation occlusion, and roughly ten pixels on target for species identification. Their
drone survey of riverine rabbits still needed ground confirmation of every sighting.
Karp~\cite{karp2020detecting} measured a 36\% detection probability for hare leverets over
109 handheld searches, with vegetation height dominating. Ulhaq et
al.~\cite{ulhaq2021automated} showed that transfer learning partly compensates for poor
spatial resolution in airborne thermal video. Mowen et al.~\cite{mowen2022improving}
classified animal poses from vehicle-mounted FLIR imagery at 82\% accuracy while noting that
a single-city dataset limits generalisation. Reviews of remotely piloted
survey~\cite{hyun2020remotely} and of monitoring methods for large
mammals~\cite{prosekov2020methods} reach the same conclusion: thermal sensors detect
homeothermic animals well, but automated species identification from thermal imagery alone
remains unreliable. McMahon et al.~\cite{mcmahon2021comparing} compared UAS thermal survey
against pellet-group counts for white-tailed deer and obtained overlapping but very wide
intervals, 4--24 deer\,km$^{-2}$. They remain the closest published comparison for our
species, and they report density rather than individual counts. Nazir and
Kaleem~\cite{nazir2021advances} note that thermal work lags the RGB domain in dataset scale
and standardisation.

Aerial and underwater counting studies show the same reliance on per-frame evaluation.
Brown et al.~\cite{brown2022automated} tested detection across ground sample distances from
0.2 to 1.0\,m\,px$^{-1}$ and found that image detail, not network capacity, sets the limit.
Infantes et al.~\cite{infantes2022automated} built an end-to-end pinniped survey from flight
planning to automatic body-size measurement. Kline et al.~\cite{kline2024integrating}
tracked zebra and giraffe from an autonomous uncrewed aerial vehicle (UAV) but analysed herd
centroids rather than individuals, and Mpouziotas et al.~\cite{mpouziotas2023automated}
reached 91.3\% AP on drone bird footage evaluated per frame only. Kilfoil et
al.~\cite{kilfoil2020using} found that both human and CNN counts of sea cucumbers
underestimate at high density. Tarling et al.~\cite{tarling2022deep} applied density
regression with uncertainty quantification to sonar imagery of fish schools, using
self-supervised pretraining to compensate for a small labelled set. Automated bird detection
has also been evaluated for airport strike management~\cite{oviedo2024automatic}.

Deep learning for camera traps was catalysed by Norouzzadeh et
al.~\cite{norouzzadeh2018automatically}, who matched crowd-sourced labels on 3.2 million
Snapshot Serengeti images. Beery et al.~\cite{beery2018recognition} then showed that models
trained at one camera location fail at unseen ones, and Kellenberger et
al.~\cite{kellenberger2018detecting} found the same failure across terrain in aerial
survey. Singh et al.~\cite{singh2020animal} documented limited transfer from natural to
man-made environments, and Beery et al.~\cite{beery2020synthetic} showed that synthetic
renders reduce error for rare classes only. Liu et al.~\cite{liu2023lote} released a
twelve-year longitudinal video dataset and confirmed that web-image pretraining transfers
only partially to camera-trap footage. Adam et al.~\cite{adam2025wildlifereid10k} designed a
similarity-aware split for a re-identification benchmark to prevent the near-duplicate
leakage that inflates reported accuracy, the same contamination our held-out protocol
(\cref{sec:heldout}) excludes. Throughout this literature the unit of analysis is the
trigger event, and individual identity is assumed unknown unless the species carries
distinctive markings.

One absence runs through all of it. These methods treat imperfect detection as a parameter
to be estimated, not as a sequence of pipeline stages whose losses can be separated. The
abundance estimators model it as a detection function, the imaging studies evaluate per
frame, and the camera-trap literature takes the trigger event as its unit. We found no
thermal or camera-based wildlife study that annotates individual animals as distinct tracks
across continuous video, and none that reports a distinct-individual count evaluated against
per-individual ground truth. That is the gap this paper addresses. Its ground truth is
itself subject to imperfect detection, since an annotated animal is one a reviewer could see
in the footage. The decomposition therefore attributes the losses of an automated pipeline
relative to what is visible in the video, not to the true number of animals present
(\cref{sec:operating}).

\subsection{Detection and Tracking at Small Scale}
\label{sec:rw:detect}

COCO defines small objects as area below $32^2$~\cite{lin2014coco}, and wildlife at range
falls well under it. Wang et al.~\cite{wang2021nwd} show that IoU is extremely sensitive to localisation error at small scale. A one-pixel
diagonal shift drops the IoU of a $6\times6$ box to 0.53 and a four-pixel shift to 0.06,
whereas a $36\times36$ box keeps 0.90 and 0.65 under the same shifts. Their
Normalized Wasserstein Distance models boxes as Gaussians and gains 6.7 AP on tiny-object
benchmarks. Zhu et al.~\cite{zhu2021tph} added a transformer prediction head and attention
to target the drone scale distribution. \Cref{sec:criterion} makes the same sensitivity exact
for our box size, and \cref{sec:choosing} shows its consequence: at 27 pixels the strict
threshold discards detections a counting task would accept, and IoU$\geq$0.50 re-ranks the
detectors relative to a counting criterion.

Both one- and two-stage families appear in our benchmark: Faster R-CNN~\cite{ren2015faster}
as the two-stage baseline, ATSS~\cite{zhang2020atss} and TOOD~\cite{feng2021tood} as
anchor-assignment refinements, RTMDet~\cite{lyu2022rtmdet} as a modern convolutional
one-stage design, and the YOLO line~\cite{redmon2016yolo,wang2024yolov9,wang2024yolov10} as
the practical default. Query-based detectors struggle on very small objects: DETR~\cite{carion2020detr}
also converges slowly, DINO~\cite{zhang2023dino} repairs the convergence, and
RT-DETR~\cite{zhao2024rtdetr} brings the family to real-time inference. Recent reviews trace the YOLO line from v1
to v12~\cite{sapkota2025yoloreview} and compare v8 through v11 architecturally, noting that
some blocks are unchanged across versions and that missing publications for several
releases complicate reproducibility~\cite{hidayatullah2025yolov8}. Our animals are a median
$29\times24$ pixels, squarely where these differences matter. \Cref{sec:detection} reports
how each family transfers under identical preprocessing and splits.

Tracking-by-detection dominates. SORT~\cite{bewley2016sort} and
DeepSORT~\cite{wojke2017deepsort} established the Kalman-plus-association template,
ByteTrack~\cite{zhang2022bytetrack} showed that associating low-confidence detections
recovers substantial recall, and BoT-SORT~\cite{aharon2022botsort} added camera-motion
compensation and an appearance term. Our camera moves at 25--90\,km\,h$^{-1}$, so global
motion compensation is not optional. We use BoT-SORT with sparse optical-flow compensation
and ablate its appearance term in \cref{sec:inversion}. Evaluation has moved toward
higher-order tracking accuracy (HOTA), which separates detection from association accuracy,
and Valmadre et al.~\cite{valmadre2021local} argue that global metrics mask local failures.
MOTRv2~\cite{zhang2023motrv2} shows that even end-to-end transformer trackers benefit from
strong detection priors. Our \reachedq/\primaryq/\countedq decomposition parallels the
detection/association split but adds a third axis, confirmation, that these metrics do not
measure. Multi-object tracking benchmarks score identity \emph{preservation}, not whether an
object should be counted at all.

\subsection{Counting, Identity and Evaluation}
\label{sec:rw:counting}

Density regression, pioneered by MCNN~\cite{zhang2016single}, estimates per-image counts
without detection, but the estimate is per frame and does not deduplicate across frames.
That suits dense scenes with hundreds of overlapping instances. Ours is the opposite regime:
sparse instances, each of which must be resolved as an individual because survey statistics
are computed per animal. Video-level counting is structurally different, since it requires
detection, association, and a rule that decides whether an object is \emph{new}. Dolokov et
al.~\cite{dolokov2023upper} exploit a known maximum animal count to reduce identity switches
in laboratory mice, but report tracking quality rather than a unique-individual count. Our
contribution is to show that in the sparse regime the detection-then-count framing is not wrong so much as \emph{mis-attributed}. The detector is near-saturated, and the accuracy a
practitioner would try to buy with a better backbone is lost downstream.

Individual identification from appearance is established for species with high-contrast
markings. Photo-identification networks reach high accuracy on coat
patterns~\cite{miele2020revisiting}, deformable transformers reach 85\% mAP on cross-camera
tiger data~\cite{li2024redeformtr}, and WildlifeReID-10k spans 33 species and 10{,}000
identities~\cite{adam2025wildlifereid10k}. All operate on high-resolution crops of
distinctively patterned animals. \Cref{sec:reid} tests whether this transfers and finds that
it does not: at a median 27-pixel animal, appearance carries less identity signal than a
two-number size descriptor, and an ImageNet backbone scores \emph{below} plain geometry. We
therefore make no individual-identification claim and do not use appearance matching in the
reported system.

Nguyen et al.~\cite{nguyen2023trustworthy} show that detection rankings fluctuate with the
IoU threshold and propose formal requirements that average precision does not satisfy. LRP
Error~\cite{oksuz2022lrp} separates localisation from classification quality and yields a
deployment threshold. Cai and Vasconcelos~\cite{cai2021cascade} document the high-quality
detection paradox, where raising the training IoU threshold causes overfitting as positives
vanish, and Boundary IoU~\cite{cheng2021boundary} shows that mask IoU over-penalises small
objects. Our finding that IoU$\geq$0.50 re-ranks seven of eleven detectors by two or more
places, and that the best AP$_{50}$ model counts fewer animals than the one we adopt,
instantiates exactly the concern Nguyen et al.\ raise. On the calibration side, Rahman et
al.~\cite{rahman2021perframe} predict per-frame mAP from internal features to monitor
deployment without ground truth. Tarling et al.~\cite{tarling2022deep} show that
uncertainty quantification improves both training and downstream biological decisions,
which is the role our calibrated per-track posterior plays in \cref{sec:calibration}.
\label{sec:othermodalities}

The closest methodological comparison is the Caltech Fish Counting
benchmark~\cite{kay2022caltech}, which poses the same problem in sonar: count distinct
individuals in continuous video, not detections per frame. Its conclusion is the opposite of
ours. Substituting a perfect detector recovers almost all of the counting error, so in that
setting detection is the binding constraint. Sonar imagery has a low signal-to-noise ratio, and the
reported failures are generalization failures of detection across deployments. A thermal
deer at 27 pixels is faint but consistently present, and a vehicle transect gives every
animal tens to hundreds of frames of visibility with strong ego-motion cues for association.
What generalizes from this paper is therefore not ``detection is solved'' but the
\emph{instrument}: a staged decomposition that reports which stage binds, so that a
practitioner in a new modality can measure rather than assume. \Cref{sec:loso} reinforces
the point, since detection recovers some ground as a bottleneck once sites change.

That framing also identifies the literature our confirmation stage should be read against.
Duplicate suppression at the detection level is mature: Soft-NMS re-scores rather than
discards overlapping boxes~\cite{bodla2017softnms}, and Confluence replaces the IoU
criterion with a proximity measure that behaves better on small
objects~\cite{shuai2021confluence}. All of it operates \emph{within} a frame. The duplicates
that limit counting here are between \emph{tracks}, separated in time, on the same animal,
and \cref{sec:nms} tests whether the frame-level tools reach that far. Calibration is the
other relevant thread. Modern neural networks are systematically
overconfident~\cite{guo2017calibration}, which is why \cref{sec:calibration} reports a
calibrated per-track posterior rather than a raw score, and why \cref{sec:loso} finds that
an absolute score threshold does not survive a change of site.

\subsection{Positioning}

\begin{table}[htbp]
\caption{\textbf{Where this work sits.} The closest comparables in thermal wildlife
survey, automated aerial counting, and animal re-identification. ``n/r'' means the study does
not report the quantity.\label{tab:related}}
\centering\footnotesize
\setlength{\tabcolsep}{3pt}
\begin{tabularx}{\textwidth}{Llllcccc}
\toprule
\textbf{Work} & \textbf{Modality} & \textbf{Platform} & \textbf{Taxa} &
\textbf{Per-ind.\ ID} & \textbf{Video} & \textbf{Count?} & \textbf{Split} \\
\midrule
Burke 2019~\cite{burke2019optimizing} & Thermal & UAV & Rabbit & No & No & No & n/r \\
Karp 2020~\cite{karp2020detecting} & Thermal & Handheld & Hare & No & No & No & n/r \\
Ulhaq 2021~\cite{ulhaq2021automated} & Thermal & Aircraft & Multiple & No & Yes & No & n/r \\
McMahon 2021~\cite{mcmahon2021comparing} & Thermal & UAV & W.\ deer & No & No & Density & n/r \\
Corcoran 2019~\cite{corcoran2019automated} & Thermal & UAV & Koala & No & No & No & n/r \\
Mowen 2022~\cite{mowen2022improving} & Thermal & Vehicle & Multiple & No & Yes & No & n/r \\
Brown 2022~\cite{brown2022automated} & RGB & UAV & Livestock & No & No & No & Yes \\
Norouzzadeh 2018~\cite{norouzzadeh2018automatically} & RGB & Camera trap & 48 spp. & No & No & No & Site \\
Beery 2018~\cite{beery2018recognition} & RGB & Camera trap & Multiple & No & No & No & Location \\
Kellenberger 2018~\cite{kellenberger2018detecting} & RGB & Aircraft & Multiple & No & No & No & Region \\
Liu 2023~\cite{liu2023lote} & RGB & Camera trap & Endangered & No & Yes & No & Domain \\
Adam 2025~\cite{adam2025wildlifereid10k} & RGB & Mixed & 33 spp. & Yes & No & No & Similarity \\
Tarling 2022~\cite{tarling2022deep} & Sonar & Fixed & Fish & No & Yes & Density & n/r \\
Dolokov 2023~\cite{dolokov2023upper} & RGB & Overhead & Mice & Yes & Yes & No & n/r \\
\midrule
\textbf{This work} & \textbf{Thermal} & \textbf{Vehicle} & \textbf{W.\ deer} &
\textbf{Yes} & \textbf{Yes} & \textbf{Yes} & \textbf{Video} \\
\bottomrule
\end{tabularx}
\end{table}

\Cref{tab:related} makes the gap specific. Its ``per-individual ID'' column asks whether
each animal is annotated as its own track rather than as anonymous per-frame boxes. The last
two columns isolate the gap: whether a count of distinct animals is reported end to end, and
whether training and evaluation are separated by video or site rather than by frame. Thermal
datasets with per-individual track identity over continuous video are essentially absent.
Existing thermal work annotates per-frame detections or static images, and the RGB video
datasets that carry individual identity are either photo-identification collections of
still crops~\cite{adam2025wildlifereid10k} or laboratory settings with a known number of
animals~\cite{dolokov2023upper}. Train--test
separation by video or site is frequently unclear in the thermal and aerial literature. Our
own measurement quantifies what that costs: evaluating on all 32 videos rather than the 13
held out inflates coverage by 15.5 points, from 66.3\% to 81.8\% (\cref{tab:leak}). Most
distinctively, no prior work isolates \emph{which stage} of a detect--track--confirm
pipeline loses animals. The literature treats detection quality as the binding constraint;
our decomposition shows that it is not here, and that each richer candidate pool we built
made the final count worse. The closest conceptual parallel is the detection/association
split within HOTA, which measures identity preservation but not whether an object should be
counted at all.
